\section{The Pre-Merits Attribution Gap}
\label{sec:gap}

A merits rule and a fact-access rule perform different work. Negligence may ask what precautions a reasonable deployer would have taken; products law may ask whether a design or warning was defective; antidiscrimination law may ask whether a screening process caused a prohibited disparity; administrative law may ask whether an agency supplied the explanation and review an entitlement requires. Each inquiry presupposes a defined act, actor, system, and time. Affected people commonly know the adverse result and the institution through which it arrived. They do not ordinarily know the version, configuration, allocation of control, or record needed to connect that result to the governing elements.

This Part separates three forms of opacity, maps the gap along actor-identity and record-availability coordinates, identifies the predicates ordinary procedure presupposes, distinguishes adjacent evidence rights, and uses the deployment-control cascade to specify the missing record.

\subsection{The Output Does Not Reveal the Deployment}

An AI dispute should begin with the deployed system rather than the generated token, score, or recommendation. The first attribution question is not who selected the final words or number. It is who constructed, supplied, authorized, observed, and could alter the deployment that produced the legally relevant effect. A model's internal computation may be difficult to explain while the surrounding control structure can leave usable evidence in contracts, release records, configuration systems, evaluation reports, incident queues, and access logs when law requires its creation and retention. That distinction makes responsibility evidence obtainable even when a compact explanation of the model's internal path is unavailable.

Three forms of opacity should remain separate. \term{Computational opacity} concerns the inability to give a human-scale account of how internal states produced a particular output. \term{System opacity} concerns interactions among the model, instructions, retrieval, memory, tools, filters, credentials, and conventional software. \term{Organizational opacity} concerns who approved the use, possessed risk information, controlled resources, or could restrict and stop operation. The first may follow from architecture. The latter two ordinarily follow from deployment design and recordkeeping.

Collapsing all three into a single ``black box'' allows the least tractable fact to obscure the facts ordinary doctrine can use. A court may be unable to reconstruct why a model ranked one token above another yet able to determine that a provider selected the model; a deployer supplied protected records and external tools; an interface suppressed uncertainty; a safety team documented a known failure mode; management rejected a delayed release; and an operator lacked time or authority to intervene. Intent and knowledge are often inferred from conduct and circumstances, product defect can sometimes be proved circumstantially, and direct organizational fault can lie in information systems, supervision, and resource allocation.\lrfootnote{See, e.g., \casecite[444--46]{usGypsum1978}; \casecite[41--43]{speller2003}; \bluecite[\S\S~7.03(1), 7.05]{restatementAgency2006}.} Model interpretability is not a condition precedent to responsibility at the deployment level.

Fragmentation makes those accessible facts difficult to attribute. Consider an employer that licenses an applicant platform. The employer may select the job, criteria, candidate pool, weight assigned to a recommendation, and final disposition. A platform vendor may select the ranking model, training process, default features, update cadence, and bias testing. A systems integrator may configure data fields and workflow triggers. The applicant encounters an employer-branded portal and a rejection, not this allocation. In \emph{Mobley v. Workday, Inc.}, for example, a complaint alleged that employer customers delegated candidate screening to Workday's AI-based tools; at the Rule 12 stage, the district court treated the alleged delegation of a traditional hiring function as sufficient to proceed on a statutory agency theory while leaving the allegations and ultimate liability unresolved.\lrfootnote{See \casecite[806--08]{mobley2024}. The decision granted in part and denied in part Workday's motion to dismiss the first amended complaint; it did not find discriminatory operation or resolve the merits.} The opinion shows that substantive law may address a divided stack once the vendor and its alleged function are visible. It does not supply every rejected applicant with the identity, version, configuration, or evaluation record needed to reach that point.

The same problem appears when the visible institution itself cannot inspect the vendor's method. A public employer may control whether and how a score affects termination while a vendor controls the scoring software and updates. A hospital may select the patient population and integrate a recommendation while a supplier controls model validation and releases. A benefits agency may own the legal decision and appeal path while contractors control the calculation tool and documentation. Control is distributed by function, not transferred wholesale to whichever entity's logo appears on the screen.

The front door follows legally answerable organizations; it does not construct an artificial defendant. Operational autonomy describes how a system produces effects, not an entity with assets, records, capacity for service, or exposure to judgment. The rule begins with a visible organization and follows records and risk-specific authority to enterprises law can reach. A future legislature may create a bounded juridical status, but that separate project is unnecessary to open the merits process.

\subsection{Two Coordinates Define the Gap}

Actor identity and record availability are independent coordinates. Treating them as one generic opacity problem obscures which legal intervention is needed. A claimant may know the relevant operator yet lack its deployment record. A claimant may possess a transaction or application record yet lack the upstream entity that supplied the operative function. A divided stack can produce both conditions at once. The front door sequences identity and record access so that each condition receives the smallest useful response.

The first configuration is an identifiable integrated operator with an adequate record. Ordinary procedure has a direct target, and the statute's principal work occurs ex ante: standardization, retention, and a quick first answer. The record may establish the active version, relevant policy, warnings, incident sequence, and stop authority without adding another merits defendant. If a filed claim proceeds, ordinary discovery can deepen the record in the familiar way.

The second configuration is an identifiable operator with a missing or fragmented internal record. The brand and legal entity are visible, but the claimant and sometimes the operator cannot reconstruct which versions, components, policies, and intervention roles governed the event. This is the production problem. It arises in an integrated consumer service when engineering, safety, product, trust-and-safety, and executive decisions live in separate systems. It also arises in a public institution that relies on a proprietary score it cannot independently verify. The ex ante duty matters most here because a post-event subpoena cannot create lineage or approval that the operator never recorded.

The third configuration is an identifiable event with an unknown operative entity. An applicant may have an employer, job, timestamp, and rejection while lacking the product feature that processed the application and the supplier responsible for it. A patient may know the hospital encounter while lacking the vendor and version embedded in a clinical workflow. This is the actor-identification problem. A closed first answer resolves provenance and custodian fields before broader access. The visible operator performs the gateway role because it chose to incorporate the system and can preserve the event even if another enterprise holds the relevant control lane.

The fourth configuration combines the two. The operative feature and supplier are hidden, and the minimum record was never created or is divided by contract. A supplier may host data while the customer owns it; an integrator may hold configuration history while the institution holds the decision; a model provider may know version lineage while the application company knows prompts, tools, and user exposure. Asking the claimant to choose a single controller before revealing this map converts organizational design into a litigation advantage. Treating every participant as a controller creates the opposite error. The front door identifies lanes first and terminates the special role of actors outside them.

The two coordinates yield the Article's sequence: map the event and deployment, preserve the bounded slice, identify control and custody lanes, remove irrelevant actors, and expand access only along lanes relevant to an unresolved legal theory. That sequence opens the merits without deciding them.

\subsection{Ordinary Procedure Opens After Four Predicates}
\label{sec:ordinary-procedure}

Federal procedure supplies substantial tools after four predicates are present: a filed action, a complaint capable of stating a legal claim, a known party or subpoena recipient, and an existing record within the reach of that target. The deployment front door supplies information needed to establish those predicates; it does not replace the procedure that follows them.

Pleading law illustrates the sequence. Rule 8 requires a short and plain statement showing entitlement to relief. \emph{Twombly} and \emph{Iqbal} require factual matter supporting a plausible claim, while Rule 11(b)(3) requires counsel to certify that factual contentions have evidentiary support or, if specifically identified, will likely have such support after a reasonable opportunity for investigation or discovery.\lrfootnote{See \statcite{frcp8}; \statcite{frcp11}; \casecite[555--59, 570]{twombly2007}; \casecite[678--79, 686]{iqbal2009}; \casecite[560--61]{rotella2000}.} That accommodation matters: a plaintiff need not possess every internal email or validation report before filing. But Rule 11 creates neither a discovery entitlement nor a substitute for Rule 8. Neither rule names an unknown supplier, links a rejection to a product feature, identifies the version that ran, or requires the operator to have recorded that fact.

Discovery becomes powerful once an action and target exist. Rule 26(b)(1) permits discovery of relevant, proportional, nonprivileged matter. Rule 34 reaches existing documents and electronically stored information in a party's possession, custody, or control. Rule 45 reaches a known nonparty that can be named and served.\lrfootnote{See \statcite{frcp26}; \statcite{frcp34and45}.} These mechanisms can test an agency theory, obtain validation evidence, and expose a divided contract. They do not identify the person to whom a subpoena should be directed or create a release approval, configuration history, or incident record that never existed.

Initial disclosure begins later and serves a narrower function. Rule 26(a)(1) generally requires the parties to identify witnesses and materials they may use to support their own claims or defenses, along with damages and insurance information. The disclosures ordinarily follow the Rule 26(f) conference.\lrfootnote{See \statcite{frcp26Initial}.} They are not an adverse-record inventory. An operator that does not expect to use a rejected mitigation, contradictory test, or internal escalation in support of its position has no initial-disclosure obligation to list the item on that basis alone. The claimant can request it in discovery after the case begins; the front-door problem is how to identify the system, holder, and record category needed to frame that request.

Court-authorized early discovery provides a real bridge when the claimant already has a unique event identifier and knows the intermediary holding the identity record. In \emph{Strike 3 Holdings, LLC v. Doe}, a rights holder tied an alleged online infringement to an IP address and sought an early subpoena to the known internet service provider capable of mapping that address to a subscriber. The D.C. Circuit held that Rule 26(d)(1) permits such a route and that Doe pleading can be appropriate when discovery is expected to reveal identity, subject to relevance, proportionality, and a realistic prospect of jurisdiction and success.\lrfootnote{See \casecite[1207, 1210--11]{strike3Doe2020}; see also \casecite[577--80]{seescandy1999}.} The case identifies the bridge's predicates as clearly as its power: a viable substantive claim, a unique incident identifier, and a known intermediary with an existing identity record. A fragmented deployment may lack the second or third predicate because the applicant sees only a branded portal, the operator does not retain the feature or run identifier, or the intermediary itself is hidden behind an integrator.

Rule 27 occupies another useful but bounded lane. It permits a person expecting to become a party to perpetuate testimony that risks being lost before an action can be filed. The petition must state the facts the petitioner expects to establish and the substance of the testimony to be preserved. Courts therefore distinguish perpetuation of known evidence from an investigation designed to learn whether a cause of action exists or whom to sue. In \emph{Workman v. United States Postal Service}, the Tenth Circuit affirmed denial of a petition that sought testimony to determine whether a postal-facility fire supported a claim; lack of the evidence needed to frame a complaint was not the inability to bring an action contemplated by Rule 27.\lrfootnote{See \statcite{frcp27}; \casecite[240--45]{workman2025}; \casecite[912]{ashCort1975}.} A deposition cannot supply an unknown model version or a deployment record that no witness or enterprise created.

Divided control creates a further limit within ordinary discovery. Rule 34 control is a legal and organizational inquiry, not a synonym for technical hosting. The Ninth Circuit has declined to treat cooperative affiliation as control when the responding entity lacked a legal right to obtain an affiliate's files on demand. The Eleventh Circuit, in a different statutory proceeding and factual record, found control where affiliated entities regularly shared responsive information and the recipient had the practical ability to obtain it.\lrfootnote{Compare \casecite[1107--08]{citricAcid1999}, with \casecite[1200--02]{sergeeva2016}. \emph{Sergeeva} arose under 28 U.S.C. \S~1782; it illustrates a practical-control approach rather than deciding ordinary domestic merits discovery.} Those outcomes properly depend on contract and practice. They also show why a visible operator should retain the minimum record locally and obtain a contractual right to retrieve defined provenance and event fields from a supplier. Otherwise the vendor may physically host the data while the customer legally controls it, leaving each able to identify the other as the proper source.

Doe pleading and later amendment preserve some claims but do not create a national limitations safety net. Rule 15(c)(1)(C) requires timely notice to the added party and knowledge that, but for a mistake concerning identity, the action would have been brought against it. \emph{Krupski} treats the prospective defendant's knowledge as central and recognizes that confusion over the roles of related, known entities can constitute a mistake. Courts can distinguish that situation from complete ignorance of an upstream actor. In \emph{Winzer}, the Fifth Circuit held that substituting named officers for Doe defendants after limitations expired did not correct a Rule 15 mistake. Rule 15(c)(1)(A) can incorporate a more forgiving state relation-back rule, but that route depends on the law governing limitations and the state's own conditions.\lrfootnote{See \statcite{frcp15}; \casecite[547--57]{krupski2010}; \casecite[470--71]{winzer2019}; \casecite[1200--02]{butler2014}.} A deployment statute can address the timing problem directly for a claim it creates or expressly governs by tolling the applicable period while a compliant identity request is pending. Amendment and relation back in a later federal action remain governed by Rule 15.

State pre-suit proceedings confirm that a narrow identity stage is administrable while revealing the absence of a uniform route. Illinois Rule 224 authorizes an independent verified action limited to identifying a potentially responsible person or entity and ordinarily ends the proceeding after identification. Texas Rule 202 permits a verified petition to investigate a potential claim when the likely benefit outweighs burden or when perpetuation may prevent a failure or delay of justice. New York authorizes disclosure to aid in bringing an action. California's perpetuation statute takes the opposite approach: it expressly excludes discovery to determine whether a claim exists or to identify parties.\lrfootnote{See \statcite{illinoisRule224}; \statcite{texasRule202}; \statcite{newYorkCPLR3102}; \statcite{californiaCCP2035}. A New York pre-action petition does not become a federal Rule 27 proceeding simply because a respondent removes it; \emph{Teamsters Local 404} treated the state petition as outside the removable civil action described by 28 U.S.C. \S~1441. See \casecite[264--69]{teamstersKing2018}.} These systems provide design elements---verification, a closed purpose, burden review, a short proceeding, and cost allocation---rather than an existing federal solution.

\subsection{Existing Evidence Rights Need a Deployment Front End}
\label{sec:adjacent-literature}

Evidence-access scholarship now reaches both sides of litigation's filing line. Stewart's counterfactual-interrogation right would let a person affected by a known adverse algorithmic decision generate behavioral evidence before filing suit. The organization would supply a bounded interface using the decision-time model version, supported by logging and retention duties.\lrfootnote{See \bluecite[1067, 1075--79]{stewartBeyondExplanation2026}.} This is the nearest pre-litigation intervention. It begins, however, with an identified adverse decision, an organization capable of furnishing the testing portal, and a model linked to the decision. The front door constructs the identity and routing predicates that a divided deployment may conceal: whether a covered feature participated, which product and version ran, which entities supplied and operated it, and who can retrieve the event record. Once Stage One supplies that map, counterfactual interrogation can operate as one bounded form of Stage Two testing.

Cen, Ismael, and Zheng address the complementary access dispute. They identify seven recurring asymmetries---models, data, documentation, logs, expertise, compute, and infrastructure---and describe a Catch-22 in which the requester needs evidence to justify access while the producing party controls the evidence needed to make that showing. Their three-part test compares the benefits and risks of requested access, the minimum access needed to reproduce elements of the cause of action, and functionally equivalent alternatives.\lrfootnote{See \bluecite[7794--95, 7801--05]{cenPrivatizationProof2026}. Their analysis includes nonparty developers, preservation, procurement, and logging. The distinction here is institutional sequence: the front door imposes a visible-operator gateway, minimum deployment record, material-control routing rule, and safe exit before a court can apply an access test to a defined requester--producer relation.} That framework helps a court decide \emph{how much} existing evidence an identified producing party should disclose. Stage One supplies the deployment, event, custodian, and control map; Stage Two can then use ordinary proportionality, protective orders, access alternatives, and the Cen framework on a claim-specific record.

Other work specifies adjacent components without supplying this claimant-facing route. García-Llorente and Olmeda propose a minimum technical-evidence package for sufficient intelligibility and auditability; procurement scholarship treats the contract file as a governance record; and organizational research shows that accountability records can change the conduct they observe.\lrfootnote{See \artcite{garciaLlorenteTransparency2026} (abstract) (proposing a minimum technical-evidence package); \bluecite{maxwellProcurementRecord2026} (abstract) (treating the procurement file as a governance record); \artcite[558--63]{chappidiAccountabilityCapture2025} (identifying behavioral effects of accountability records).} Those contributions inform what a useful record may contain and how its costs should be controlled. The front door supplies the legal sequence that makes the record reachable by an affected person without treating the procurement contract or technical package as the source of the entitlement.

A 2026 systematic study of 559 federal opinions likewise finds that courts principally use preexisting doctrines and that litigated harms reflect the injuries those doctrines recognize.\lrfootnote{See \bluecite[1--2, 8--9]{yuVisibleCourt2026}. The study codes federal opinions in which AI appears in the parties' contentions; it does not estimate the universe of AI incidents or adjudicate the validity of any claim.} That evidence supports the Article's ordinary-law premise. The front door adds a positive-law routing and record-formation layer so existing doctrine can operate on a known deployment.

\subsection{The Cascade Specifies the Missing Record}
\label{sec:cascade}

The deployment-control cascade is a disciplined way to define those predicates. Adapted from systems-safety analysis, it represents an AI deployment as linked controllers, control actions, process models, feedback channels, and intervention capacity. Systems-theoretic methods treat loss as a control problem that can arise from an absent constraint, unsafe action, inadequate feedback, or a controller's mistaken model of the operating environment.\lrfootnote{See \bookcite[75--76, 87--90]{levesonEngineering2011} (developing the systems-theoretic premise); \bluecite[14--15]{levesonThomasSTPA2018} (specifying the STPA method); \artcite[2:5--2:6, 2:14--2:18]{rismaniPlaneCrashes2023} (applying safety-engineering frameworks to responsible machine learning).} The cascade uses that vocabulary only to specify contestable evidence. Governing law assigns duty, fault, causation, defenses, and remedy.

Four functional levels provide a reusable map for the integrated and divided deployments addressed here:

\begin{enumerate}[leftmargin=2.2em]
\item \term{authority and resources}: the actors who set objectives, risk tolerance, staffing, compute, data access, and release conditions;
\item \term{model and component supply}: the actors who develop or select models, document limits, conduct evaluations, issue updates, and supply interfaces or weights;
\item \term{deployment and harness operation}: the actors who select instructions, retrieval, memory, filters, routing, tools, credentials, eligible users, monitoring, and escalation; and
\item \term{the affected environment}: operators, users, data subjects, third-party systems, and institutional conditions that supply inputs and experience effects.
\end{enumerate}

These are roles rather than four firms. One enterprise can occupy the first three levels in an integrated service. An API deployment may divide them among a model vendor, application developer, enterprise customer, and professional user. A public agency can set the legal objective and use of a score while a contractor supplies and updates the model. The map prevents a corporate label or contract caption from substituting for actual function.

\begin{figure}[htbp]
\centering
\begin{tikzpicture}[
  node distance=0.72cm,
  box/.style={draw, rounded corners=2pt, align=center, text width=0.65\textwidth, minimum height=0.7cm, inner sep=5pt},
  down/.style={-{Latex[length=2mm]}, thick},
  up/.style={-{Latex[length=2mm]}, thick, dashed}
]
\node[box] (authority) {Authority and resources\\\footnotesize objective, risk tolerance, budget, staffing, release approval};
\node[box, below=of authority] (provider) {Model and component supply\\\footnotesize model selection, evaluation, documentation, updates};
\node[box, below=of provider] (deployer) {Deployment and harness operation\\\footnotesize instructions, retrieval, memory, tools, safeguards, monitoring, intervention};
\node[box, below=of deployer] (environment) {Affected environment\\\footnotesize operators, users, data subjects, institutions, and experienced effects};

\draw[down] ([xshift=2.5cm]authority.south) -- node[right, font=\scriptsize] {constraints and resources} ([xshift=2.5cm]provider.north);
\draw[down] ([xshift=2.5cm]provider.south) -- node[right, font=\scriptsize] {components and updates} ([xshift=2.5cm]deployer.north);
\draw[down] ([xshift=2.5cm]deployer.south) -- node[right, font=\scriptsize] {outputs and interventions} ([xshift=2.5cm]environment.north);

\draw[up] ([xshift=-2.5cm]environment.north) -- node[left, font=\scriptsize] {complaints and incidents} ([xshift=-2.5cm]deployer.south);
\draw[up] ([xshift=-2.5cm]deployer.north) -- node[left, font=\scriptsize] {evaluations and limits} ([xshift=-2.5cm]provider.south);
\draw[up] ([xshift=-2.5cm]provider.north) -- node[left, font=\scriptsize] {risk reports and requests} ([xshift=-2.5cm]authority.south);
\end{tikzpicture}
\caption{The deployment-control cascade as a record specification. Solid arrows carry authority, resources, components, and operational actions toward the affected environment. Dashed arrows carry the feedback required to update each controller's understanding. Restriction, rollback, credential revocation, and shutdown must be recorded where they can be exercised.}
\label{fig:cascade}
\end{figure}

At each level, six questions specify the record. What asserted risk was within scope? Which action could create, expose, detect, or reduce it? What did the controller understand about the system and environment? Which feedback could correct that understanding? What practical authority and resources existed at the relevant time? Who could restrict, roll back, suspend, or terminate operation? Part~IV translates those questions into a minimum control sheet. The answers are factual: governing law decides whether an approval, alert, credential, intervention, or omission establishes any merits element, and the same record may rebut the claim.

This mapping yields the front door's object. Actor-identification asks which enterprise occupied a risk-specific control lane. Deployment-record production asks for the bounded materials that establish or rebut that proposition and then permit ordinary doctrine to operate. Part II turns those two functions into a statutory trigger, a first-answer obligation, and a safe exit.
